\subsection{Problem Formulation and Setup}
\label{sec:taxonomy}

% \definecolor{ColorExtrapolation}{RGB}{240, 245, 255}
% \definecolor{ColorImputation}{RGB}{240, 255, 240}
% \definecolor{ColorAugmentation}{RGB}{255, 245, 235}
% \definecolor{BorderGray}{RGB}{180, 180, 180}

\definecolor{ColorExtrapolation}{RGB}{240, 245, 255}
\definecolor{ColorImputation}{RGB}{240, 255, 240}
\definecolor{ColorAugmentation}{RGB}{255, 245, 235}
\definecolor{BorderGray}{RGB}{180, 180, 180}

To establish a rigorous conceptual foundation, Figure~\ref{fig:taxonomy} presents a comprehensive two-level hierarchical taxonomy that structures the entire financial time-series generation landscape across task formulations and algorithmic families.
Let $\mathbf{X} = \{x_1, x_2, \dots, x_T\}$ denote a multivariate financial time series, where each $x_t \in \mathbb{R}^d$ represents a vector of financial variables (e.g., asset prices, returns, trading volumes, or limit order book states) at time step $t$. Financial time series are inherently stochastic, heavily influenced by non-stationary market regimes, and exhibit complex inter-variable dependencies.
The core objective of Financial Time-Series Generation (FTSG) is to learn a probabilistic generative model $P_\theta$ capable of producing synthetic sequences $\hat{\mathbf{X}}$ that capture both the marginal distributions and the temporal dynamics of the true data-generating process $P_{data}$.

\subsection{Task-Level Taxonomy}
As illustrated in Figure~\ref{fig:tasks}, FTSG comprises three distinct problem formulations based on generation context.

Financial time-series extrapolation is a conditional, forward-looking generation task. Given a historical prefix $\mathbf{x}_{1:t}$, the goal is to generate future trajectories $\mathbf{x}_{t+1:T}$. The predictive distribution is modeled as $P(\mathbf{x}_{t+1:T} \mid \mathbf{x}_{1:t}; \theta) = \prod_{i=t+1}^{T} P(x_i \mid \mathbf{x}_{1:i-1}; \theta)$. This encompasses probabilistic forecasting, where the output is a set of plausible future paths rather than a single point estimate.

Financial time-series imputation is a bidirectional, closed-format generation task. Let $\mathbf{M} \in \{0, 1\}^T$ be a binary mask indicating missing values. Given an incomplete sequence with missing elements $\mathbf{x}_{miss} = \mathbf{x} \odot \mathbf{M}$ and observed elements $\mathbf{x}_{obs} = \mathbf{x} \odot (1-\mathbf{M})$, the objective is to reconstruct the missing data while preserving local and global statistical coherence: $P(\mathbf{x}_{miss} \mid \mathbf{x}_{obs}; \theta)$.

Synthetic sequence generation, or data synthesis, is an unconditional or conditionally guided generation task aimed at creating entirely new data manifolds. The model learns to map a latent noise variable $\mathbf{z} \sim \mathcal{N}(0, \mathbf{I})$ (and optionally a condition $c$ such as market regime or volatility target) to a realistic sequence: $\hat{\mathbf{X}} = G_{\theta}(\mathbf{z}, c)$, aiming to minimize the divergence between the generated distribution $P_g$ and the real distribution $P_{data}$.

\subsection{Technique-Level Taxonomy}
To map the rapid evolution of algorithmic approaches, the underlying generative techniques can be grouped into six families.

Statistical and probabilistic methods include classical approaches such as ARIMA~\citep{box1970} and GARCH~\citep{bollerslev1986garch}, which model conditional mean and variance based on rigid mathematical assumptions.

Discriminative deep sequence models include recurrent neural networks (RNNs)~\citep{elman1990finding} and long short-term memory (LSTM) networks~\citep{hochreiter1997lstm} optimized via pointwise loss functions.

Adversarial generative models (GANs)~\citep{goodfellow2014gan} use generator-discriminator dynamics, including the Wasserstein GAN with gradient penalty (WGAN-GP)~\citep{arjovsky2017wgan,gulrajani2017wgangp}, to implicitly learn the data distribution.

Diffusion and score-based models~\citep{ho2020ddpm,song2019generative,song2021score} learn to reverse a gradual noising process, offering stable training and high-fidelity distribution matching.

Foundation models and LLM-based models are recent adaptations of large pre-trained Transformers~\citep{vaswani2017attention,brown2020language} to the time-series domain, represented by models such as Chronos~\citep{ansari2024chronos}, OneFitsAll~\citep{zhou2023onefitsall}, and Time-LLM~\citep{jin2024timellm}.

Controllable and multimodal generative models, including TimeWeaver~\citep{narasimhan2024timeweaver}, T2S~\citep{T2S}, and dual-sourced text--time diffusion models~\citep{lozano2023dual}, are designed to generate sequences conditioned on textual instructions, macroeconomic indicators, or specific volatility regimes.

The detailed classification of FTSG methods, organized by task and technique, is summarized in Figure ~\ref{fig:taxonomy} for a clear overview.

\subsection{Mapping Tasks to Techniques}
Table \ref{tab:mapping} provides a unified mapping of these tasks to the corresponding techniques, highlighting the evolution from deterministic prediction to flexible, scenario-based generative paradigms. Specifically, for the extrapolation task, the landscape spans from deep sequence architectures such as SGP-LSTM~\citep{SGP-LSTM1} that forecast prices and returns using deterministic metrics like MSE, to advanced diffusion and foundation models such as TimeGrad~\citep{rasul2021timegrad}, Chronos~\citep{ansari2024chronos}, and FinCast~\citep{zhu2025fincast}. These newer paradigms handle multivariate and multi-asset series, shifting the evaluation focus toward proper probabilistic scoring rules such as the continuous ranked probability score (CRPS)~\citep{gneiting2007strictly} and zero-shot capabilities. In the realm of imputation, missing fundamentals and prices are primarily addressed using tensor methods such as RegTensor~\citep{zhou2023regtensor} and diffusion models such as CSDI~\citep{tashiro2021csdi} and TimeDiT~\citep{cao2024timedit}, where performance is gauged by reconstruction error and temporal coherence. Finally, for data synthesis, generative frameworks such as QuantGAN~\citep{wiese2020} and CoFinDiff~\citep{tanaka2025cofindiff} are employed to simulate realistic returns, market scenarios, and complex limit order flows. Consequently, the evaluation criteria for synthesis shift entirely away from point-wise accuracy, prioritizing instead the preservation of stylized facts, distributional fidelity, and the practical downstream utility of the generated data.

